\chapter{Επαλήθευση και Ακεραιότητα Ισχυρισμών}
\label{chap:validation}

\section{Η ανάγκη για ξεχωριστή επαλήθευση}
\label{sec:validation-scope}
Στον κύβο του Rubik, η παραγωγή μιας ακολουθίας κινήσεων δεν αρκεί για επιστημονικό ισχυρισμό. Πρέπει να ελέγξουμε αν η ακολουθία λύνει πράγματι την κατάσταση, αν το μήκος της μετρήθηκε στη συμφωνημένη μετρική, αν η αρχική κατάσταση είναι φυσικά έγκυρη και αν υπάρχει απόδειξη ότι δεν υπάρχει συντομότερη λύση. Τα ερωτήματα αυτά είναι συναφή αλλά όχι ταυτόσημα, και η παρούσα εργασία τα διαχωρίζει σε διαφορετικά επίπεδα ελέγχου.

Η επαλήθευση λύσης ελέγχει μόνο την εγκυρότητα της ακολουθίας. Η επαλήθευση ακριβούς απόστασης ελέγχει επιπλέον ότι το μήκος είναι το ελάχιστο δυνατό, όταν υπάρχει πλήρης αναζήτηση κατά πλάτος (\eng{BFS}) ή ολοκληρωμένη επαναληπτικά εμβαθυνόμενη αναζήτηση A* (\eng{iterative deepening A*}, IDA*) με παραδεκτή ευρετική. Η επαλήθευση των παραγόμενων τεκμηρίων (\eng{generated artifacts}) ελέγχει ότι οι πίνακες και τα δεδομένα των σχημάτων προέρχονται από αποθηκευμένα αποτελέσματα. Τέλος, η επαλήθευση του κειμένου ελέγχει αν οι διατυπώσεις της εργασίας αντιστοιχούν στα παραπάνω τεκμήρια. Αν λείπει ένα από αυτά τα επίπεδα, μπορεί να καταλήξουμε σε ένα πρόγραμμα τεχνικά σωστό αλλά ακαδημαϊκά ασθενές. Το παρόν κεφάλαιο ορίζει τι εγγυάται κάθε επίπεδο ελέγχου και πού σταματούν οι εγγυήσεις αυτές. Η μηχανική υλοποίηση των αντίστοιχων εργαλείων περιγράφεται στο Κεφάλαιο~\ref{chap:implementation} και η θέση τους στην αρχιτεκτονική του συστήματος στο Κεφάλαιο~\ref{chap:design}.

\section{Ανεξάρτητος επαληθευτής λύσεων}
\label{sec:validation-verifier}
Ο επαληθευτής λύσεων (\eng{verifier}) ξεκινά από την αποθηκευμένη αρχική κατάσταση και εφαρμόζει την προτεινόμενη λύση στο μοντέλο κυβιδίων (\eng{cubie model}), χωρίς να χρησιμοποιεί καμία εσωτερική κατάσταση του επιλυτή. Αυτό είναι κρίσιμο: ένας επιλυτής μπορεί να κρατά δικά του εσωτερικά αντικείμενα αναζήτησης ή δικές του ενδιάμεσες καταστάσεις. Αν ο επαληθευτής εμπιστευόταν αυτά τα αντικείμενα, θα μπορούσε να αναπαραγάγει το ίδιο σφάλμα. Αντίθετα, η επανεκτέλεση από την αποθηκευμένη κατάσταση ψηφίδων (\eng{facelet state}) και την ακολουθία κινήσεων ελέγχει τη λύση ανεξάρτητα, εκ των έξω. Η ανεξαρτησία αυτή ισχύει ως προς την εσωτερική κατάσταση και τη λογική αναζήτησης του επιλυτή. Η επανεκτέλεση χρησιμοποιεί το ίδιο μοντέλο κινήσεων της υλοποίησης, οπότε ένα συστηματικό σφάλμα στο ίδιο το μοντέλο του κύβου δεν θα εντοπιζόταν από τον έλεγχο αυτόν, αλλά από τις διασταυρώσεις με εξωτερικές υλοποιήσεις που περιγράφονται στην Ενότητα~\ref{sec:validation-exact}.

Το αποτέλεσμα της επαλήθευσης είναι δυαδικό, αλλά συνοδεύεται από επεξηγηματικό μήνυμα και από το μετρημένο πλήθος κινήσεων. Αν η λύση δεν λύνει την κατάσταση, η γραμμή αποτελέσματος δεν επιτρέπεται να φέρει \code{verified=true}. Αν το δηλωμένο \code{solution_length} δεν συμφωνεί με το πλήθος κινήσεων που προκύπτει από τη συντακτική ανάλυση της ακολουθίας, ο έλεγχος αποτελεσμάτων καταγράφει σφάλμα. Έτσι αποφεύγονται σιωπηρές ασυνέπειες, όπου η ακολουθία είναι σωστή αλλά το μήκος αναφέρεται λανθασμένα.

Η ανεξαρτησία του επαληθευτή έχει και επιπτώσεις στη διατύπωση του κειμένου. Όταν το κείμενο γράφει «επαληθευμένη λύση», εννοεί ότι η ακολουθία πέρασε από αυτή τη διαδικασία· δεν εννοεί ότι ο επιλυτής δήλωσε ο ίδιος επιτυχία. Η διαφορά είναι μικρή στη διατύπωση, αλλά μεγάλη στην αξιοπιστία.

\section{Έλεγχος των ακριβών αποτελεσμάτων}
\label{sec:validation-exact}
Η ετικέτα \code{exact} (αποδεδειγμένα ελάχιστο μήκος) είναι αυστηρότερη από την ένδειξη επαλήθευσης. Για να χαρακτηριστεί ένα αποτέλεσμα ακριβές, πρέπει να υπάρχει αποδεικτική διαδικασία που αποκλείει κάθε συντομότερη λύση μέσα στην ίδια μετρική. Στις περιπτώσεις μικρού βάθους αυτό επιτυγχάνεται με εξαντλητική BFS. Στην IDA* επιτυγχάνεται όταν η αναζήτηση με παραδεκτή ευρετική ολοκληρωθεί και επιστρέψει λύση πριν εξαντληθούν τα όρια χρόνου και κόμβων. Αν η IDA* διακοπεί από όριο, δεν αποδίδεται ετικέτα \code{exact}, ακόμη και αν το αποδεδειγμένο κάτω φράγμα είναι μεγάλο.

Ο έλεγχος αποτελεσμάτων διασταυρώνει τις γραμμές \code{exact} με τη διαθέσιμη απόσταση BFS, όπου αυτή υπάρχει. Ο έλεγχος αυτός δεν καλύπτει κάθε δυνατή ακριβή περίπτωση, καλύπτει όμως τις ρηχές περιπτώσεις όπου υπάρχει ανεξάρτητη εξαντλητική απόδειξη. Αν ένας επιλυτής δήλωνε ακριβές μήκος διαφορετικό από την απόσταση BFS, ο έλεγχος θα το κατέγραφε ως σφάλμα.

Η αυστηρότητα αυτή αποτρέπει ένα συχνό λάθος: το να θεωρηθεί ακριβής η πρώτη λύση που βρέθηκε από αναζήτηση περιορισμένου βάθους ή από μια πρακτική μέθοδο φάσεων. Η πρώτη λύση μιας αναζήτησης κατά βάθος (\eng{DFS}) ή μιας πρακτικής μεθόδου φάσεων δεν είναι κατ' ανάγκην ελάχιστη· μόνο η πλήρης εξάντληση όλων των μικρότερων ορίων, ή μια ισοδύναμη απόδειξη, παρέχει τέτοια εγγύηση. Η ίδια αυστηρότητα ισχύει και για τις αναζητήσεις με αποκοπή κινήσεων ρίζας. Ο εγγενής βέλτιστος επιλυτής επιστρέφει τότε τη διακριτή ετικέτα \code{exact_under_root_mask}. Η αναβάθμιση σε πλήρες \code{exact}, χωρίς περιορισμό, γίνεται μόνο από το περίβλημα. Προϋπόθεση είναι η μάσκα να προέρχεται από τον πιστοποιημένο σταθεροποιητή συμμετρίας 48 στοιχείων που αντιστοιχεί στην αρχική κατάσταση· η αναβάθμιση καταγράφεται ρητά στο πεδίο \code{notes}.

Η απόσταση 20 του superflip είναι το μοναδικό ακριβές αποτέλεσμα για τον κύβο 3×3×3 που στηρίζεται σε χωριστή εξαντλητική απόδειξη κάτω φράγματος. Η απόδειξη αυτή βασίζεται σε ευρετική, της οποίας η παραδεκτότητα επικυρώθηκε έναντι του ακατέργαστου πίνακα BFS της φάσης~1. Η κατασκευή της απόδειξης περιγράφεται στο Κεφάλαιο~\ref{chap:implementation} και τα αριθμητικά της τεκμήρια παρουσιάζονται στο Κεφάλαιο~\ref{chap:experiments}. Τα βαθύτερα ακριβή αποτελέσματα (βάθη 15 έως 25) δεν συνοδεύονται από αντίστοιχη χωριστή απόδειξη. Στηρίζονται στην ολοκληρωμένη παραδεκτή αναζήτηση της μηχανής επίλυσης (\eng{backend}) που τα παρήγαγε. Όπου υπάρχουν πολλαπλές γραμμές για την ίδια κατάσταση, διασταυρώνονται και με ανεξάρτητες μηχανές. Στο κύριο σώμα μετρήσεων (\eng{benchmark}, Κεφάλαιο~\ref{chap:experiments}), οι γραμμές \code{exact} διασταυρώνονται με εξαντλητική BFS μέχρι την απόσταση 5. Οι βαθύτερες αποδεικνύονται με ολοκληρωμένη IDA* εντός ορίου βάθους 10, με βαθύτερη αποδεδειγμένη απόσταση την 9. Επιπλέον, η πιστοποίηση βελτιστότητας που εκτελείται μέσα στη σουίτα δοκιμών φτάνει την απόσταση 6 για την υλοποίηση καθαρής Python, με σημείο αναφοράς το μαντείο (\eng{oracle}). Φτάνει επίσης την απόσταση 11 μέσω του αμοιβαίου ελέγχου, κατά τον οποίο ο επιλυτής Korf και η μηχανή απόδειξης βελτιστότητας της σχεδίασης δύο φάσεων πρέπει να αποδείξουν ανεξάρτητα το ίδιο μήκος.

Για τις βαθύτερες αυτές γραμμές προέκυψε ένα πρόβλημα με την προέλευση των τεκμηρίων. Εκ των υστέρων διαπιστώθηκε ότι ο πίνακας αποκοπής (\eng{pruning table}) \code{h48h7}, ο οποίος στηρίζει τη μηχανή μαντείου H48, παράχθηκε με πολλά νήματα και η παραγωγή αυτή δεν ήταν ντετερμινιστική. Όταν ο πίνακας αναπαράχθηκε δοκιμαστικά από την αρχή, το άθροισμα ελέγχου βγήκε διαφορετικό. Η σύγκριση byte προς byte έδειξε την αιτία. Λόγω συνθήκης συναγωνισμού (\eng{race condition}) μεταξύ των νημάτων, περίπου 650 από τα περίπου \(7.6\times10^{9}\) στοιχεία των 4 bit (\eng{nibbles}) του διατηρούμενου πίνακα φέρουν μη ενημερωμένες τιμές, υψηλότερες από τις πραγματικές. Οι τιμές αυτές εμφανίζονται μόνο στα δευτερεύοντα (εφεδρικά) ελάχιστα κάθε γραμμής, όχι στην κύρια τιμή αποκοπής. Ο έλεγχος της κατανομής ως προς την αναμενόμενη (κανονική) κατανομή πέρασε με επιτυχία και πιστοποιεί την κύρια περιοχή αποκοπής ως ακριβή. Η γεννήτρια τροποποιήθηκε στη συνέχεια ώστε να ενημερώνει τα ελάχιστα με ατομικές πράξεις σύγκρισης και ανταλλαγής (\eng{compare-and-swap}). Ακόμη όμως και ο νέος πίνακας της δοκιμαστικής αναπαραγωγής, παραγμένος με τον διορθωμένο κώδικα, έχασε λίγες σπάνιες ενημερώσεις ελαχίστων: έξι μη ενημερωμένα στοιχεία των 4 bit. Επομένως, ο αρχικός πίνακας δεν μπορεί να αναπαραχθεί ταυτόσημα (byte προς byte), και ο διατηρούμενος πίνακας δεν αντικαταστάθηκε. Έξι τεκμήρια ενσωματώνουν το αρχικό άθροισμα ελέγχου: πέντε αρχεία πιστοποίησης του μαντείου και ένα αρχείο μεταδεδομένων του πίνακα. Παραμένουν ως έχουν, ως συνειδητή και τεκμηριωμένη απόφαση και όχι ως τεκμήρια παραγμένα εκ νέου. Πρόκειται για περιορισμό προέλευσης των συγκεκριμένων τεκμηρίων, όχι για ελάττωμα ορθότητας κάποιου αναφερόμενου ακριβούς αποτελέσματος.

Επειδή τιμές αποκοπής ξεπερασμένες και υπερεκτιμημένες (προς τα πάνω) θα μπορούσαν, κατ' αρχήν, να αποκόψουν μια βέλτιστη λύση, η εγκυρότητα των βαθιών γραμμών \code{exact} δεν στηρίζεται στον πίνακα αυτόν. Οι 11 αυτές γραμμές είχαν ως μοναδικό τεκμήριο τον πίνακα \code{h48h7}. Επιβεβαιώθηκαν ανεξάρτητα από το εξωτερικό Nissy 2.0.8, με τους δικούς του δημόσιους πίνακες~\cite{trontoNissyCoreH48}: 11/11 συμφωνίες ως προς το μήκος και επιτυχής αναπαραγωγή κάθε λύσης μέχρι τη λυμένη κατάσταση στο μοντέλο κύβου της παρούσας υλοποίησης (\path{results/processed/h48only_rows_independent_reverify_seed_2026_thesis.json}). Έτσι κάθε βαθιά ακριβής γραμμή διαθέτει τουλάχιστον μία ανεξάρτητη διασταύρωση.

\section{Έλεγχος πινάκων και ευρετικών}
\label{sec:validation-tables}
Οι πίνακες αποκοπής χρησιμοποιούνται ως παραδεκτά κάτω φράγματα. Η παραδεκτότητα στηρίζεται στο ότι κάθε πίνακας αποθηκεύει πραγματική απόσταση σε προβολή του πλήρους προβλήματος. Για να έχει νόημα αυτός ο ισχυρισμός, οι πίνακες πρέπει να έχουν παραχθεί από σωστές μεταβάσεις. Οι δοκιμές ελέγχουν ότι οι πίνακες κινήσεων συμφωνούν με τις άμεσες κινήσεις στο μοντέλο κυβιδίων και ότι οι κύκλοι κωδικοποίησης και αποκωδικοποίησης των συντεταγμένων είναι συνεπείς.

Ο έλεγχος αυτός δεν αποδεικνύει ότι το φράγμα είναι ισχυρό· αποδεικνύει ότι το φράγμα αντιστοιχεί στην προβολή που δηλώνει. Η ισχύς του φράγματος φαίνεται πειραματικά από το πόσο επιταχύνει την αναζήτηση, ενώ η αδυναμία του αποκαλύπτεται από τις γραμμές με ετικέτα \code{timeout} (υπέρβαση του ορίου χρόνου). Η εργασία χρειάζεται και τα δύο στοιχεία: μαθηματική βεβαιότητα ότι η ευρετική δεν υπερεκτιμά, αλλά και πειραματική παραδοχή ότι δεν αρκεί πάντα.

\section{Έλεγχος του κύβου 2×2×2}
\label{sec:validation-pocket}
Η πλήρης μελέτη του κύβου 2×2×2 (\eng{Pocket Cube}) έχει ξεχωριστή επαλήθευση, επειδή είναι το μόνο σημείο όπου η εργασία παράγει πλήρη κατανομή αποστάσεων για ολόκληρο τον χώρο καταστάσεων. Ο επαληθευτής ελέγχει τέσσερα πράγματα: ότι η κατανομή είναι πλήρης· ότι το πλήθος των καταστάσεων που προσπέλασε η BFS ισούται με το θεωρητικό μέγεθος του κανονικοποιημένου χώρου (Κεφάλαιο~\ref{chap:experiments})· ότι το άθροισμα των κλάσεων της κατανομής ισούται με το ίδιο πλήθος· και ότι οι αντιπροσωπευτικές λύσεις είναι \code{exact} και επαληθευμένες. Οι έλεγχοι αυτοί καλύπτουν τόσο το συνολικό όσο και το επιμέρους επίπεδο, την κάθε περίπτωση ξεχωριστά.

Πέρα από τους ελέγχους εσωτερικής συνέπειας, ο επαληθευτής συγκρίνει ολόκληρη την κατανομή αποστάσεων ανά βάθος με τις κανονικές τιμές του κανονικοποιημένου μοντέλου 2×2×2. Η μέγιστη απόσταση που προκύπτει από την εξαντλητική αναζήτηση είναι 11, με ακριβώς 2644 αντιποδικές καταστάσεις (καταστάσεις σε μέγιστη απόσταση) σε αυτή την απόσταση· η τιμή 11 ταυτίζεται με τον δημοσιευμένο αριθμό του Θεού (\eng{God's number}) για τον κύβο 2×2×2 στη μετρική μισών στροφών (\eng{half-turn metric}, HTM)~\cite{randelshoferPocketCube}. Ο έλεγχος αυτός αποτρέπει μια λεπτή κατηγορία σφάλματος: μια εσφαλμένη κωδικοποίηση συντεταγμένων ή ένας λανθασμένος πίνακας μετάβασης θα μπορούσε να παραγάγει κατανομή φαινομενικά σωστή: πλήρη, με σωστό άθροισμα, που περνά τους προηγούμενους ελέγχους, αλλά στην ουσία λανθασμένη. Επιπλέον, ένας ανεξάρτητος έλεγχος διασταύρωσης επαληθεύει ότι οι πίνακες μετάβασης των συντεταγμένων συμφωνούν με το μοντέλο κυβιδίων. Οι πίνακες αυτοί προκύπτουν από ανάλυση σε μετάθεση και προσανατολισμό, και ο έλεγχος καλύπτει κάθε μετάβαση των ρηχών προσβάσιμων καταστάσεων. Έτσι τεκμηριώνεται ότι η ανάλυση αναπαράγει πιστά τις πραγματικές κινήσεις του κύβου και δεν είναι απλώς αυτοσυνεπής.

Η πληρότητα εδώ έχει σαφή όρια: αφορά το κανονικοποιημένο μοντέλο 2×2×2 με σταθερό σημείο αναφοράς, όχι τον πλήρη κύβο 3×3×3, και η διατύπωση αυτή διατηρείται σε όλο το κείμενο. Η αξία του ελέγχου είναι ότι δείχνει σε τι συνίσταται μια πραγματικά πλήρης απόδειξη σε εφικτό χώρο. Η περίπτωση 2×2×2 δεν χρησιμοποιείται για να καλύψει κενά του 3×3×3· χρησιμοποιείται ως καθαρό παράδειγμα σωστής πειραματικής μεθοδολογίας.

\section{Έλεγχος των παραγόμενων πινάκων του κειμένου}
\label{sec:validation-artifacts}
Το κείμενο της εργασίας περιέχει παραγόμενους πίνακες LaTeX και αρχεία δεδομένων για σχήματα. Τα αρχεία αυτά δεν αντιμετωπίζονται ως περιεχόμενο γραμμένο από τον συγγραφέα, σε αντιδιαστολή με το αυτόματα παραγόμενο: είναι παράγωγα των αποθηκευμένων αρχείων αποτελεσμάτων και αναγεννώνται όταν αλλάζουν τα αποτελέσματα. Το εργαλείο ελέγχου της εργασίας εντοπίζει μη ενημερωμένα παράγωγα συγκρίνοντας τους χρόνους τροποποίησης των επεξεργασμένων πηγών με εκείνους των παραγόμενων αρχείων (Κεφάλαιο~\ref{chap:implementation}). Αν ένας πίνακας είναι παλαιότερος από τη σύνοψη από την οποία προέρχεται, υπάρχει κίνδυνος το τελικό έγγραφο να εμφανίζει παλαιούς αριθμούς.

Ο έλεγχος αυτός είναι πρακτικός αλλά όχι τέλειος: δεν αποδεικνύει ότι το περιεχόμενο ενός πίνακα είναι σημασιολογικά σωστό. Γι' αυτό συνδυάζεται με τον επαληθευτή αποτελεσμάτων και με τις δοκιμές. Εντοπίζει όμως ένα σημαντικό είδος αστοχίας: αλλαγή των μετρήσεων χωρίς αντίστοιχη αναγέννηση των πινάκων του κειμένου. Σε πειραματική εργασία αυτό το είδος αστοχίας είναι συχνό και επικίνδυνο.

\section{Δείκτες κινδύνου ισχυρισμών}
\label{sec:validation-markers}
Το εργαλείο ελέγχου της εργασίας εντοπίζει επίσης επίμαχες λέξεις που εγκυμονούν κίνδυνο για τη διατύπωση ισχυρισμών, όπως «βέλτιστος», «ακριβής», «πλήρης» και «αυθαίρετος». Οι λέξεις αυτές δεν είναι απαγορευμένες· απαιτούν όμως τεκμηρίωση. Για παράδειγμα, ο τίτλος της εργασίας περιέχει τη λέξη «βέλτιστη» επειδή αυτός είναι ο τίτλος του θέματος. Επίσης, η εισαγωγή δηλώνει ρητά ότι η εργασία δεν παράγει βέλτιστες λύσεις για αυθαίρετες καταστάσεις. Και οι δύο αυτές διατυπώσεις είναι αποδεκτές. Αν όμως μια πρόταση χαρακτήριζε τον εγγενή επιλυτή Kociemba πλήρη βέλτιστο επιλυτή, θα ήταν εσφαλμένη.

Κάθε τέτοιο σημείο του κειμένου εξετάστηκε ένα προς ένα από τον συγγραφέα και αντιστοιχίστηκε στο είδος τεκμηρίου που το στηρίζει: βιβλιογραφική παραπομπή, δοκιμή, μέτρηση, απόδειξη ή ρητά διατυπωμένος περιορισμός. Όπου δεν υπήρχε τεκμήριο, η διατύπωση αναθεωρήθηκε. Το εργαλείο υποδεικνύει πού πρέπει να εστιάσει κανείς· η κρίση για το αν ένας ισχυρισμός στηρίζεται επαρκώς παραμένει ευθύνη του συγγραφέα.

\section{Σχέση επαλήθευσης και περιορισμών}
\label{sec:validation-limitations}
Οι περιορισμοί, όταν καταγράφονται με ακρίβεια, δεν αποτελούν ένδειξη αποτυχίας· αποτελούν μέρος της επαλήθευσης. Αν ένας εγγενής επιλυτής εμφανίζει εξαντλήσεις ορίου χρόνου, η καταγραφή τους είναι η ορθή συμπεριφορά. Αν δεν υπάρχει βέλτιστος επιλυτής με εγγυημένο πρακτικό χρόνο για κάθε αυθαίρετη κατάσταση του κύβου 3×3×3, το κείμενο οφείλει να το δηλώνει ρητά. Η απόκρυψη ενός περιορισμού θα ήταν σοβαρότερο πρόβλημα από τον ίδιο τον περιορισμό.

Διακρίνονται πάντως δύο είδη περιορισμών. Οι περιορισμοί εμβέλειας οριοθετούν τι αποδεικνύεται. Δεν υπάρχει καθολική εγγύηση πρακτικού χρόνου για τον πλήρη κύβο 3×3×3· αυτό το όριο μετριάζεται από τα ακριβή αποτελέσματα με φραγμένα όρια αναζήτησης και από την πλήρη μελέτη του 2×2×2, και δηλώνεται ρητά ως όριο της συνεισφοράς. Αντίθετα, ένα ελάττωμα ορθότητας δεν θα ήταν αποδεκτό σε καμία εμβέλεια — για παράδειγμα, μια λύση που δεν επαληθεύεται ή ένα ακριβές μήκος που αναιρείται από ανεξάρτητη απόδειξη. Ολόκληρος ο μηχανισμός επαλήθευσης του παρόντος κεφαλαίου υπάρχει για να αποκλείει αυτό το δεύτερο είδος.

\section{Η αρχή της επαληθευσιμότητας}
\label{sec:validation-principle}
Η γενική αρχή της εργασίας είναι ότι κάθε ισχυρισμός πρέπει να μπορεί να αναχθεί σε αποθηκευμένο τεκμήριο. Αν μια πρόταση μιλά για πίνακα, υπάρχει παραγόμενος πίνακας. Αν μιλά για πλήθος γραμμών, υπάρχει αρχείο αποτελεσμάτων JSON. Αν μιλά για ακριβή λύση, υπάρχει αναζήτηση που την αποδεικνύει. Αν μιλά για περιορισμό, υπάρχει τεκμηριωμένη καταγραφή του. Η αυστηρή αυτή προσέγγιση κάνει την παραγωγή των αποτελεσμάτων πιο αργή, αλλά την ανάγνωσή τους πιο αξιόπιστη.

Η ίδια αρχή ισχύει για κάθε μελλοντική επέκταση. Δεν αρκεί να προστεθεί ταχύτερος επιλυτής· πρέπει να προστεθούν τρόπος επαλήθευσης, μετρήσεις, σχήμα αποτελεσμάτων και διατύπωση στο κείμενο εξίσου ακριβής. Μόνο τότε μια τεχνική βελτίωση γίνεται ακαδημαϊκή συνεισφορά.
